\chapter{Perché un MILP con errore \texorpdfstring{$L_1$}{L1}}
\label{ch:perche-l1}

\section{Il vincolo imposto dal solver}

HiGHS risolve problemi di programmazione lineare (LP) e lineare mista intera (MILP), ma
\textbf{non accetta termini quadratici}. Non esiste inoltre una funzione obiettivo simbolica come
in altri solver: l'obiettivo è un \textbf{vettore di coefficienti} $c$, uno per ciascuna variabile
del problema, e il solver minimizza il prodotto scalare
\begin{equation}
  \min_{v} \; c^{\mathsf{T}} v .
\end{equation}

Questa è una restrizione forte e determina tutto ciò che segue: \textbf{ogni penalità del modello
deve essere esprimibile come un coefficiente costante moltiplicato per una variabile}. Non si
possono penalizzare prodotti di variabili, né funzioni non lineari di esse. Quando nel seguito
compare una penalità apparentemente complessa --- per esempio proporzionale a quanto un ciclo
sfora la durata massima --- la sua realizzazione consiste sempre nell'introdurre una variabile
ausiliaria e assegnarle un coefficiente lineare.

\section{Dall'errore quadratico all'errore assoluto}

La scelta naturale per misurare la qualità della ricostruzione sarebbe l'errore quadratico
\begin{equation}
  \sum_{t} \bigl( \hat{y}(t) - y(t) \bigr)^2 ,
  \qquad
  \hat{y}(t) = \sum_i P_i x_i(t),
\end{equation}
che però è quadratico e quindi inammissibile. Si adotta allora l'errore assoluto ($L_1$)
\begin{equation}
  \sum_{t} \bigl| \hat{y}(t) - y(t) \bigr| ,
\end{equation}
che pur non essendo lineare può essere \textbf{linearizzato esattamente}.

\subsection{La linearizzazione}

Si introducono per ogni istante due variabili continue non negative, $\ep_t$ ed $\en_t$, e si
impone l'uguaglianza
\begin{equation}
  \hat{y}(t) \;-\; \ep_t \;+\; \en_t \;=\; y(t),
  \qquad \ep_t \ge 0, \quad \en_t \ge 0 ,
  \label{eq:linearizzazione}
\end{equation}
minimizzando poi la somma $\ep_t + \en_t$ con coefficiente unitario su entrambe.

Il meccanismo è il seguente. La~\eqref{eq:linearizzazione} si riscrive come
$\ep_t - \en_t = \hat{y}(t) - y(t)$: le due variabili devono avere differenza pari al residuo, ma
la loro \emph{somma} non è vincolata dall'alto. Poiché entrambe compaiono nell'obiettivo con
costo positivo, il solver non ha alcun interesse a gonfiarle: all'ottimo almeno una delle due vale
sempre zero, e di conseguenza
\begin{equation}
  \ep_t + \en_t \;=\; \bigl| \hat{y}(t) - y(t) \bigr| .
\end{equation}
Non serve alcun vincolo aggiuntivo per imporre questa proprietà --- discende dal segno dei
coefficienti. In pratica $\ep_t$ misura di quanto il modello \emph{sovrastima} il consumo reale in
quell'istante ed $\en_t$ di quanto lo \emph{sottostima}; una sola delle due è attiva per volta.

\begin{modelbox}{Termine di ricostruzione}
\begin{equation*}
  \sum_{t \in \Vset} \bigl( \ep_t + \en_t \bigr)
  \qquad\text{soggetto a}\qquad
  \hat{y}(t) - \ep_t + \en_t = y(t) \quad \forall t \in \Vset
\end{equation*}
\end{modelbox}

\section{Conseguenza sulla scala delle penalità}

Questa scelta ha un effetto diretto sulla taratura di tutto il modello, ed è il punto più
importante da tenere presente leggendo i capitoli seguenti.

L'errore di ricostruzione è misurato in \textbf{watt}. Perché una penalità sia commensurabile con
esso --- cioè perché abbia senso sommarla nella stessa funzione obiettivo --- deve avere la stessa
dimensione fisica. È per questa ragione che \textbf{tutte le penalità del modello sono scalate per
la potenza tipica $p_i$ dell'elettrodomestico a cui si riferiscono}: i coefficienti $\lambda$ sono
adimensionali e il prodotto $\lambda \cdot p_i$ restituisce watt.

La conseguenza pratica è un'interpretazione uniforme e immediata di ogni coefficiente:
\begin{center}
\begin{minipage}{0.88\textwidth}
\itshape
un coefficiente pari a 1 significa che violare quella dichiarazione del questionario una volta
costa esattamente quanto lasciare $p_i$ watt non spiegati per la durata di uno slot.
\end{minipage}
\end{center}

Questo è il \emph{punto di pareggio}, e su di esso si basa tutta la taratura discussa nel
capitolo~\ref{ch:parametri}. Sotto 1 la penalità è più debole dell'evidenza del segnale e il
modello preferisce spiegare i dati; sopra 1 prevale la dichiarazione dell'utente.
